\DeclareEditionOverlay{CM-C05-U001}{%
\begin{EditionUnitCard}{Un mapa de trabajo}
Predice qué información aportará cada cálculo antes de hacerlo.
\RegisteredBookFigure{FIG-C05-001}{fig:c05:function-study-workflow}{Flujo razonado de un estudio cualitativo.}{Diagrama enlazado desde dominio y simetrías hasta límites, derivadas, representación y comprobación, con un retorno cuando aparece una incoherencia.}{\FunctionStudyWorkflowVisual}
\begin{EditionCheckpointBox}{Comprueba}
Explica qué decisión gráfica depende de cada etapa.
\end{EditionCheckpointBox}
\end{EditionUnitCard}}

\DeclareEditionOverlay{CM-C05-U002}{%
\begin{EditionUnitCard}{Simetría con dominio}
Antes de comparar f(-x), comprueba que x y -x sean admisibles.
\begin{EditionWarningBox}{Atención}
La fórmula sola no basta si el dominio no es simétrico.
\end{EditionWarningBox}
\end{EditionUnitCard}}

\DeclareEditionOverlay{CM-C05-U003}{%
\begin{EditionUnitCard}{Repetición exacta}
La periodicidad es una identidad compatible con el dominio, no una semejanza visual.

\end{EditionUnitCard}}

\DeclareEditionOverlay{CM-C05-U004}{%
\begin{EditionUnitCard}{Intersecciones admisibles}
Filtra las soluciones por el dominio antes de colocarlas en la gráfica.
\begin{EditionCheckpointBox}{Comprueba}
Comprueba por separado eje horizontal y eje vertical.
\end{EditionCheckpointBox}
\end{EditionUnitCard}}

\DeclareEditionOverlay{CM-C05-U005}{%
\begin{EditionUnitCard}{Cada asíntota tiene su límite}
Nombra primero hacia dónde se aproxima la variable.
\begin{EditionCheckpointBox}{Comprueba}
Vincula vertical, horizontal u oblicua con el límite que la justifica.
\end{EditionCheckpointBox}
\end{EditionUnitCard}}

\DeclareEditionOverlay{CM-C05-U006}{%
\begin{EditionUnitCard}{Haz que las tablas conversen}
La información de primera y segunda derivada debe producir una gráfica coherente.
\begin{EditionWarningBox}{Atención}
Si las conclusiones chocan, revisa signos e intervalos antes de dibujar.
\end{EditionWarningBox}
\end{EditionUnitCard}}

\DeclareEditionOverlay{CM-C05-U007}{%
\begin{EditionUnitCard}{Síntesis completa}
Construye la gráfica por capas y justifica cada rasgo.
\begin{EditionCheckpointBox}{Comprueba}
Antes de ver el resultado final, bosqueja una predicción con dominio y asíntotas.
\end{EditionCheckpointBox}
\end{EditionUnitCard}}

\DeclareEditionOverlay{CM-C05-U008}{%
\begin{EditionUnitCard}{Comprobar no es demostrar}
Usa software después de predecir y para detectar incoherencias.
\begin{EditionWarningBox}{Atención}
Una ventana puede ocultar comportamiento; conserva la justificación analítica.
\end{EditionWarningBox}
\end{EditionUnitCard}}

